\section{Estructura de les peticions i respostes JSON}
\label{annex:json}

Aquest annex mostra una versió abreujada, però fidel, del contracte real de l'endpoint \texttt{/analyze}. Els nivells de risc interns utilitzats pel \textit{backend} són \texttt{BAJO}, \texttt{MEDIO} i \texttt{ALTO}, i les accions són \texttt{ALLOW}, \texttt{REVIEW} i \texttt{BLOCK}. La DApp i el Snap en mostren els equivalents en català, però el contracte conserva aquests valors per compatibilitat.

\subsection{Petició enviada pel Snap}

\begin{lstlisting}[language=json]
{
  "chainId": "eip155:11155111",
  "from": "0x852712de...45AE881aa",
  "to": "0xA32D1b3B...c672a363",
  "value": "0x0",
  "data": "0x095ea7b3...0001",
  "origin": "https://sepolia.etherscan.io"
}
\end{lstlisting}

El camp \texttt{data} conté la \textit{calldata}; els primers quatre bytes formen el selector. En aquest exemple, \texttt{0x095ea7b3} correspon a \texttt{approve(address,uint256)}.

\subsection{Resposta representativa del backend}

\begin{lstlisting}[language=json]
{
  "success": true,
  "risk": "BAJO",
  "risk_score": 20,
  "issues": [
    "S'ha detectat una aprovaci\u00f3 de tokens",
    "Quantitat aprovada: 1",
    "L'aprovaci\u00f3 no \u00e9s il\u00b7limitada"
  ],
  "verdict": {
    "risk": "BAJO",
    "risk_score": 20,
    "recommended_action": "ALLOW",
    "source": "deterministic_base",
    "reason": "No hi ha hagut motius suficients per alterar el veredicte base"
  },
  "decoded": {
    "method": "approve",
    "signature": "approve(address,uint256)",
    "selector": "0x095ea7b3",
    "spender": "0x3333333333333333333333333333333333333333",
    "amount": "1",
    "is_infinite_approval": false
  },
  "ai_review": {
    "ai_risk_hint": "BAJO",
    "confidence": "baja",
    "ai_flags": [],
    "reviewer_summary": "Sense observacions addicionals de la IA.",
    "raw_response": null
  },
  "final_verdict": {
    "risk_level": "BAJO",
    "source": "deterministic_base",
    "reason": "No hi ha hagut motius suficients per alterar el veredicte base"
  },
  "analysis_id": 48,
  "performance": {
    "ai_review_ms": 29,
    "explanation_ms": 10,
    "persistence_ms": 38,
    "total_backend_ms": 93
  },
  "evaluation": null,
  "timestamp": "2026-07-18T16:35:41.289Z"
}
\end{lstlisting}

La resposta diferencia el veredicte determinista, la revisió complementària, el veredicte final i les mesures de rendiment. \texttt{raw\_response} conserva la sortida original d'Ollama quan existeix; si el servei no està disponible, queda a \texttt{null} i el sistema utilitza una explicació local segura. Els indicis del camp \texttt{findings} són cadenes de text, no objectes amb camps \texttt{code} o \texttt{severity}.

\subsection{Persistència i detall}

\texttt{analysis\_id} identifica el registre d'\texttt{analysis\_history}. El detall recuperat des de \texttt{/analysis-history/:id} afegeix els camps complets de transacció normalitzada, senyals de memòria local, coincidències amb adreces conegudes i fets semàntics. El \textit{hash} d'una transacció confirmada no forma part de la resposta prèvia a la confirmació, perquè encara no existeix en el moment de l'anàlisi.
